% Figure: how the uncertainty signal scales trade size and protects capital
\begin{figure}[htbp]
\centering
\begin{tikzpicture}[
  box/.style={draw, rounded corners=3pt, minimum height=1.4cm,
              text width=#1, align=center, font=\small},
  arr/.style={-{Stealth[length=3mm]}, thick},
]
  \node[box=2.7cm, fill=blue!10] (unc)
    {Higher\\uncertainty $u_t$};
  \node[box=2.9cm, fill=purple!10, right=0.6cm of unc] (size)
    {Smaller\\trade size\\$\times \max(1-u_t,\,s_{\min})$};
  \node[box=2.7cm, fill=yellow!12, right=0.6cm of size] (exp)
    {Lower market\\exposure};
  \node[box=2.7cm, fill=green!12, right=0.6cm of exp] (dd)
    {Shallower\\drawdown};

  \draw[arr] (unc.east)  -- (size.west);
  \draw[arr] (size.east) -- (exp.west);
  \draw[arr] (exp.east)  -- (dd.west);

  % Hard guard note feeding into the exposure step
  \node[box=5cm, fill=red!8, below=1.0cm of size, xshift=1.7cm] (guard)
    {Hard guard: if $u_t \geq \tau$, block new buys entirely};
  \draw[arr] (guard.north) -- ++(0,0.3) -| (exp.south);

  \node[below=0.6cm of guard, text width=12cm, align=center, font=\footnotesize\itshape]
    {When the forecaster is unsure, the agent trades smaller (and stops buying
     altogether past the threshold $\tau$), so it carries less risk into turbulent
     markets and its losses stay shallow.};
\end{tikzpicture}
\caption{How the uncertainty signal scales trade size: the more unsure the forecaster
is, the smaller the agent trades --- and above the threshold $\tau$ it stops buying
altogether --- which lowers exposure in turbulent markets and keeps drawdowns shallow.}
\label{fig:uncertainty_trade_scaling}
\end{figure}
